\documentclass[11pt,letterpaper]{article}
\usepackage[margin=1in]{geometry}
\usepackage{amsmath,amssymb,bm}
\usepackage{booktabs}
\usepackage{enumitem}
\usepackage{xcolor}
\usepackage[colorlinks=true,linkcolor=blue!70!black,urlcolor=blue!70!black,citecolor=blue!70!black]{hyperref}
\usepackage{fancyhdr}
\usepackage{titlesec}
\usepackage{tocloft}
\usepackage{algorithm}
\usepackage{algpseudocode}
\usepackage{mathtools}

\pagestyle{fancy}
\fancyhf{}
\rhead{SO(3) Attitude Filter --- A Lie Group Tutorial}
\lhead{}
\rfoot{Page \thepage}
\setlength{\headheight}{14pt}

\setlength{\parskip}{5pt}
\setlength{\parindent}{0pt}

\titlespacing*{\section}{0pt}{14pt}{6pt}
\titlespacing*{\subsection}{0pt}{10pt}{4pt}
\titlespacing*{\subsubsection}{0pt}{8pt}{3pt}

% Convenience commands
\newcommand{\R}{\mathbb{R}}
\newcommand{\SO}{\mathrm{SO}(3)}
\newcommand{\so}{\mathfrak{so}(3)}
\newcommand{\skew}[1]{[#1]_{\times}}
\newcommand{\hatmap}[1]{#1^{\wedge}}
\newcommand{\veemap}[1]{#1^{\vee}}
\newcommand{\expmap}{\mathrm{Exp}}
\newcommand{\logmap}{\mathrm{Log}}
\newcommand{\eye}{\mathbf{I}}

\title{\textbf{Building an Attitude Filter on SO(3):\\
A Lie Group Extended Kalman Filter from Scratch}\\[8pt]
\large For Engineers Who Already Know Kalman Filtering}
\author{}
\date{February 2026}

\begin{document}
\maketitle
\thispagestyle{fancy}

\tableofcontents
\newpage

%======================================================================
\section{Introduction: Why This Document?}
%======================================================================

If you already know how to build a Kalman filter, you know how to
build a Lie group filter.  You just need a few new pieces of
vocabulary.  This document develops those pieces one at a time and
then assembles them into a working \textbf{Extended Kalman Filter
on $\SO$}---an attitude filter that estimates the orientation of a
rigid body from gyroscope and accelerometer measurements.

The goal is economy: only the Lie theory that is directly needed,
motivated by the estimation problem at every step.  No root systems,
no Dynkin diagrams, no abstract manifold theory.  Just matrices,
calculus, and filtering.

\subsection*{What You Should Already Know}

\begin{itemize}[nosep]
  \item The standard discrete-time Extended Kalman Filter (EKF):
    predict/update cycle, state vector, covariance matrix,
    Jacobians, Kalman gain.
  \item Basic linear algebra: matrix multiplication, transpose,
    inverse, eigenvalues.
  \item Rotation matrices and what they represent (a coordinate
    frame's orientation in 3D space).
\end{itemize}

\subsection*{What You Will Learn}

\begin{enumerate}[nosep]
  \item Why Euler angles and raw quaternions cause problems in a
    Kalman filter.
  \item The five Lie theory tools you need: the group $\SO$, the
    algebra $\so$, the exponential map, the logarithmic map, and
    the right Jacobian.
  \item How to build the ``error-state'' EKF on $\SO$ step by step.
  \item A complete algorithm you can implement in MATLAB or Python.
\end{enumerate}

\newpage
%======================================================================
\section{The Problem with Euler Angles\\(and Why We Need Lie Groups)}
%======================================================================

Suppose you want to estimate the orientation of a drone, a
satellite, or a robot link.  Your first instinct is to put three
Euler angles $(\phi, \theta, \psi)$ into the state vector and run
a standard EKF.  This works---until it doesn't.

\subsection{Gimbal Lock}

At $\theta = \pm 90^{\circ}$, the Euler angle parameterisation is
singular: two of the three axes align, a degree of freedom is lost,
and the Jacobians blow up.  This is not a numerical issue---it is
a \emph{topological} fact.  No three-parameter representation of
3D rotation can be globally non-singular.

\subsection{Quaternion Normalisation}

Quaternions avoid gimbal lock, but they have four components subject
to the constraint $\|q\| = 1$.  If you put four quaternion
components into the state vector, the covariance matrix is
$4 \times 4$ but the system has only 3 rotational degrees of
freedom.  The covariance is always rank-deficient, which causes
numerical problems.

\subsection{The Root Cause}

The space of 3D rotations is \emph{not} $\R^{3}$.  It is a curved,
compact, three-dimensional manifold.  Any attempt to treat it as a
flat vector space (by using Euler angles, Rodrigues parameters,
or unconstrained quaternions as the state) will eventually produce
artefacts.

The solution is to work with the rotation group directly---and to
put the \emph{error state} (which is small) in a local vector
space where the EKF machinery applies naturally.  This is the Lie
group approach.

\newpage
%======================================================================
\section{The Five Lie Theory Tools You Need}
%======================================================================

This section introduces exactly the Lie theory required for the
filter.  Nothing more.

%----------------------------------------------------------------------
\subsection{Tool 1: The Lie Group $\SO$}
%----------------------------------------------------------------------

\begin{center}
\fbox{\parbox{0.9\textwidth}{
\textbf{Definition.} $\SO$ is the set of all $3 \times 3$ rotation
matrices:
\[
  \SO = \bigl\{ R \in \R^{3 \times 3} \;\big|\;
  R^{\mathsf{T}} R = \eye, \;\; \det(R) = +1 \bigr\}.
\]
It is a \emph{group} under matrix multiplication.}}
\end{center}

\textbf{What this means for you:}
\begin{itemize}[nosep]
  \item If $R_1$ and $R_2$ are rotations, so is $R_1 R_2$
    (closure).
  \item Every rotation $R$ has an inverse $R^{-1} = R^{\mathsf{T}}$
    (inversion is free---just transpose).
  \item The identity rotation is $\eye_{3 \times 3}$.
  \item Rotations compose by matrix multiplication, not by adding
    angles.
\end{itemize}

\textbf{Key point:}  The ``state'' of an attitude filter is an
element $R \in \SO$, not a vector in $\R^{3}$.  We never flatten
it into three numbers.  Instead, $R$ is always a $3\times 3$
orthogonal matrix.

%----------------------------------------------------------------------
\subsection{Tool 2: The Lie Algebra $\so$ (the Tangent Space)}
%----------------------------------------------------------------------

\begin{center}
\fbox{\parbox{0.9\textwidth}{
\textbf{Definition.} The Lie algebra $\so$ is the set of all
$3 \times 3$ skew-symmetric matrices:
\[
  \so = \bigl\{ \Omega \in \R^{3 \times 3} \;\big|\;
  \Omega^{\mathsf{T}} = -\Omega \bigr\}.
\]
Every element of $\so$ is determined by a 3-vector
$\bm{\phi} = (\phi_1, \phi_2, \phi_3)^{\mathsf{T}} \in \R^{3}$.}}
\end{center}

\begin{center}
\fbox{\parbox{0.9\textwidth}{
\textbf{Hat and vee maps.}  The \textbf{hat map}
$(\cdot)^{\wedge}$ converts a 3-vector to a skew-symmetric matrix:
\[
  \bm{\phi}^{\wedge}
  \;\coloneqq\;
  \begin{pmatrix}
    0 & -\phi_3 & \phi_2 \\
    \phi_3 & 0 & -\phi_1 \\
    -\phi_2 & \phi_1 & 0
  \end{pmatrix}
  \in \so.
\]
The \textbf{vee map} $(\cdot)^{\vee}$ is the inverse: it extracts
the 3-vector from a skew-symmetric matrix.  We also write
$\skew{\bm{\phi}}$ for $\bm{\phi}^{\wedge}$ when emphasising
the cross-product interpretation:
$\bm{\phi}^{\wedge}\,\mathbf{v} = \bm{\phi} \times \mathbf{v}$.}}
\end{center}

\textbf{What this means for you:}  The Lie algebra is a
\emph{vector space}---it is $\R^{3}$ in disguise.  This is where
the error state, the covariance, and the Kalman update all live.
The algebra is the ``local flat map'' of the curved rotation space.

\textbf{Physical interpretation:}  A vector
$\bm{\phi} \in \R^{3}$ in the Lie algebra represents a small
rotation.  Its direction is the rotation axis; its magnitude
$\|\bm{\phi}\|$ is the rotation angle in radians.

%----------------------------------------------------------------------
\subsection{Tool 3: The Exponential Map
  ($\R^{3} \to \SO$)}
%----------------------------------------------------------------------

\begin{center}
\fbox{\parbox{0.9\textwidth}{
\textbf{Definition.}  The exponential map takes a Lie algebra
element (a rotation vector) and produces a rotation matrix:
\[
  \expmap(\bm{\phi})
  = \exp\!\bigl(\bm{\phi}^{\wedge}\bigr)
  = \eye
    + \frac{\sin\theta}{\theta}\,\bm{\phi}^{\wedge}
    + \frac{1 - \cos\theta}{\theta^{2}}\,
      (\bm{\phi}^{\wedge})^{2},
\]
where $\theta = \|\bm{\phi}\|$.  This is the \textbf{Rodrigues
formula}.}}
\end{center}

\textbf{What this means for you:}  The exponential map is the
bridge from the vector space (where the EKF lives) to the group
(where the state lives).  When the Kalman update gives you a
correction vector $\delta\bm{\phi} \in \R^{3}$, you apply it to
the current rotation estimate via:
\[
  R_{\text{new}} = \expmap(\delta\bm{\phi}) \; R_{\text{old}}.
\]
This is the ``reset'' or ``injection'' step.

\textbf{Small-angle limit:}  When $\|\bm{\phi}\|$ is small,
\[
  \expmap(\bm{\phi}) \;\approx\; \eye + \bm{\phi}^{\wedge},
\]
a familiar first-order approximation.

%----------------------------------------------------------------------
\subsection{Tool 4: The Logarithmic Map
  ($\SO \to \R^{3}$)}
%----------------------------------------------------------------------

\begin{center}
\fbox{\parbox{0.9\textwidth}{
\textbf{Definition.}  The logarithmic map is the inverse of the
exponential map.  Given a rotation matrix $R$:
\begin{align*}
  \theta &= \arccos\!\left(\frac{\mathrm{tr}(R) - 1}{2}\right), \\
  \logmap(R) &= \frac{\theta}{2\sin\theta}
    \bigl(R - R^{\mathsf{T}}\bigr)^{\vee}.
\end{align*}
The result is a 3-vector: the axis-angle representation.}}
\end{center}

\textbf{What this means for you:}  The logarithmic map lets you
compute the ``difference'' between two rotations as a vector:
\[
  \bm{\phi}_{\text{error}}
  = \logmap\!\bigl(R_{\text{true}}\, R_{\text{est}}^{\mathsf{T}}\bigr).
\]
This error vector lives in $\R^{3}$ and can be used in least-squares
cost functions, residual computations, or filter innovations.

%----------------------------------------------------------------------
\subsection{Tool 5: The Right Jacobian of $\SO$}
%----------------------------------------------------------------------

When we propagate the covariance through the exponential map, we
need to know how a small perturbation in the Lie algebra maps to a
perturbation in the group.  This is captured by the \textbf{right
Jacobian} $\mathbf{J}_{r}(\bm{\phi})$.

\begin{center}
\fbox{\parbox{0.9\textwidth}{
\textbf{Definition.}  For $\bm{\phi} \in \R^{3}$ with
$\theta = \|\bm{\phi}\|$:
\[
  \mathbf{J}_{r}(\bm{\phi})
  = \eye
    - \frac{1 - \cos\theta}{\theta^{2}}\,\bm{\phi}^{\wedge}
    + \frac{\theta - \sin\theta}{\theta^{3}}\,
      (\bm{\phi}^{\wedge})^{2}.
\]
For small $\bm{\phi}$:
\[
  \mathbf{J}_{r}(\bm{\phi})
  \;\approx\; \eye - \tfrac{1}{2}\,\bm{\phi}^{\wedge}.
\]}}
\end{center}

\textbf{What this means for you:}  The right Jacobian plays the
role of the ``state transition Jacobian'' for the orientation error.
In the EKF covariance propagation, it appears wherever the standard
EKF would have a partial derivative of the state transition with
respect to the state.

\textbf{Useful identity:}
\[
  \mathbf{J}_{r}^{-1}(\bm{\phi})
  = \eye
    + \tfrac{1}{2}\,\bm{\phi}^{\wedge}
    + \biggl(\frac{1}{\theta^{2}}
    - \frac{1 + \cos\theta}{2\theta\sin\theta}\biggr)
    (\bm{\phi}^{\wedge})^{2}.
\]

\textbf{Small-angle approximation:}  In practice, the time step
$\Delta t$ in an attitude filter is small (e.g., 0.01\,s), so
$\bm{\phi} = \bm{\omega}\,\Delta t$ is small and
$\mathbf{J}_{r} \approx \eye$ is often a perfectly adequate
approximation.

\subsection{Summary: The Five Tools at a Glance}

\begin{center}
\begin{tabular}{lll}
\toprule
\textbf{Tool} & \textbf{What it does} & \textbf{Filter role} \\
\midrule
$\SO$ (group) & Set of rotation matrices & The state space \\
$\so \cong \R^{3}$ (algebra) & Skew-symmetric matrices / 3-vectors
  & Error-state space \\
$\expmap: \R^{3} \to \SO$ & Vector $\to$ rotation (Rodrigues)
  & Apply correction \\
$\logmap: \SO \to \R^{3}$ & Rotation $\to$ vector
  & Compute error/innovation \\
$\mathbf{J}_{r}$ (right Jacobian) & Sensitivity of $\expmap$
  & Covariance propagation \\
\bottomrule
\end{tabular}
\end{center}

\newpage
%======================================================================
\section{The Estimation Problem}
%======================================================================

We want to estimate the orientation of a rigid body (e.g., a drone
or satellite) using:
\begin{itemize}[nosep]
  \item A \textbf{gyroscope} that measures angular velocity
    $\bm{\omega}_m \in \R^{3}$ in the body frame, corrupted by
    bias and noise.
  \item An \textbf{accelerometer} that measures the gravity vector
    $\mathbf{a}_m \in \R^{3}$ in the body frame, corrupted by
    noise.  (We assume the body is not accelerating, or that
    accelerations are small compared to gravity.)
\end{itemize}

\subsection{State}

The filter state is:
\begin{align*}
  R &\in \SO && \text{(orientation: body frame $\to$ world frame)}, \\
  \mathbf{b} &\in \R^{3} && \text{(gyroscope bias)}.
\end{align*}

\textbf{Important:}  $R$ is \emph{not} in the state vector---it is
a matrix that we propagate on the group.  The state vector for the
EKF contains only the \emph{error state}
$\delta\mathbf{x} \in \R^{6}$, which lives in the Lie algebra (and
a Euclidean space for the bias).

\subsection{Error-State Definition}

We define the error state as:
\[
  \boxed{
  \delta\mathbf{x}
  = \begin{pmatrix} \delta\bm{\phi} \\ \delta\mathbf{b} \end{pmatrix}
  \in \R^{6},
  }
\]
where:
\begin{itemize}[nosep]
  \item $\delta\bm{\phi} \in \R^{3}$ is the \textbf{orientation
    error} in the Lie algebra.  It relates the true and estimated
    rotations by:
    \[
      R_{\text{true}}
      = \expmap(\delta\bm{\phi})\; R_{\text{est}}.
    \]
    When the estimate is perfect, $\delta\bm{\phi} = \mathbf{0}$
    and $\expmap(\mathbf{0}) = \eye$.
  \item $\delta\mathbf{b} = \mathbf{b}_{\text{true}} -
    \mathbf{b}_{\text{est}} \in \R^{3}$ is the bias error
    (ordinary Euclidean subtraction).
\end{itemize}

\textbf{Key insight:}  The error state is always a vector in
$\R^{6}$.  We can define a $6 \times 6$ covariance matrix $P$ for
it, compute Jacobians, and apply the standard EKF equations---all
without any singularities, because $\delta\bm{\phi}$ is always
small when the filter is working.

\subsection{Dynamics}

The continuous-time kinematics of a rotation matrix driven by
angular velocity $\bm{\omega}$ (in the body frame) are:
\[
  \dot{R} = R\,\bm{\omega}^{\wedge}.
\]

The gyroscope measures:
\[
  \bm{\omega}_m = \bm{\omega} + \mathbf{b} + \mathbf{n}_{\omega},
\]
where $\mathbf{b}$ is a slowly varying bias and
$\mathbf{n}_{\omega} \sim \mathcal{N}(\mathbf{0},\,
\sigma_{\omega}^{2}\,\eye)$ is white noise.  The bias evolves as a
random walk:
\[
  \dot{\mathbf{b}} = \mathbf{n}_{b},
  \qquad
  \mathbf{n}_{b} \sim \mathcal{N}(\mathbf{0},\,
  \sigma_{b}^{2}\,\eye).
\]

\subsection{Why Can't Gyros Alone Estimate Attitude?}

At first glance, gyroscopes seem sufficient: they measure angular
velocity, and integrating angular velocity gives attitude.  But
gyro-only attitude estimation suffers from \emph{unbounded drift},
and this drift is fundamental---not a matter of sensor quality.

\textbf{The drift mechanism.}  The gyroscope measures
$\bm{\omega}_m = \bm{\omega} + \mathbf{b} + \mathbf{n}_{\omega}$.
To propagate attitude, we integrate $\bm{\omega}_m - \hat{\mathbf{b}}$,
where $\hat{\mathbf{b}}$ is our bias estimate.  Two error sources
cause the attitude estimate to drift without bound:

\begin{enumerate}[nosep]
  \item \textbf{Bias error integrates to a ramp.}  If the bias
    estimate $\hat{\mathbf{b}}$ is wrong by
    $\delta\mathbf{b} = \mathbf{b} - \hat{\mathbf{b}}$, then the
    attitude error grows as
    $\delta\bm{\phi}(t) \approx \delta\mathbf{b}\,t$---a ramp.
    After 10 minutes with a $0.01$\,deg/s bias error, attitude
    has drifted by $6$ degrees.  Worse, the bias itself
    random-walks, so the rate of drift is unpredictable.
  \item \textbf{White noise integrates to a random walk.}  Even with
    perfect bias compensation, the white noise
    $\mathbf{n}_{\omega}$ integrates to a random walk in attitude.
    The standard deviation grows as
    $\sigma_{\phi}(t) = \sigma_{\omega}\sqrt{t}$---unbounded.
\end{enumerate}

\textbf{The observability problem.}  With gyros alone, the filter
has no way to distinguish the true attitude from the drifted
estimate.  The bias $\mathbf{b}$ is \emph{unobservable}: there is
no measurement that constrains it, so the bias covariance grows
without bound and drags the attitude covariance with it.

\textbf{What the accelerometer provides.}  An accelerometer
measures the gravity vector in the body frame.  Gravity is a
\emph{known, fixed reference direction} in the world frame
($\mathbf{g} = (0, 0, -g)^{\mathsf{T}}$).  By comparing the
predicted gravity direction $R^{\mathsf{T}}\mathbf{g}$ with the
measured direction $\mathbf{a}_m$, the filter obtains an
\emph{absolute} orientation reference that does not drift.  This
makes two of the three attitude axes (roll and pitch)
\emph{observable}, bounding their errors.

\textbf{The yaw gap.}  Gravity is a single vector, so it
constrains the tilt of the body relative to the gravity direction
(roll and pitch) but says nothing about rotation \emph{around}
the gravity vector (yaw/heading).  Yaw remains unobservable
from accelerometer alone.  To bound yaw drift, a heading
reference is needed---typically a magnetometer (measuring Earth's
magnetic field) or GPS-derived velocity.

\textbf{Summary:}  Gyros provide high-bandwidth, short-term
attitude tracking but drift without bound.  Accelerometers provide
a low-bandwidth, absolute gravity reference that prevents roll and
pitch drift.  The filter fuses both: gyros for dynamics (prediction),
accelerometer for correction (update).  This complementary structure
is why essentially every practical attitude estimation system uses
both sensors.

\subsection{Measurement Model}

The accelerometer measures the gravity vector expressed in the body
frame:
\[
  \mathbf{a}_m
  = R^{\mathsf{T}}\,\mathbf{g} + \mathbf{n}_{a},
  \qquad
  \mathbf{n}_{a} \sim \mathcal{N}(\mathbf{0},\,
  \sigma_{a}^{2}\,\eye),
\]
where $\mathbf{g} = (0, 0, -g)^{\mathsf{T}}$ is the gravity
vector in the world frame.

\newpage
%======================================================================
\section{The Filter: Step by Step}
%======================================================================

The filter follows the standard EKF predict/update cycle, but with
one crucial difference: the state $R$ is propagated on the group,
while the covariance and Kalman gain operate on the error state in
$\R^{6}$.

%----------------------------------------------------------------------
\subsection{Step 0: Initialisation}
%----------------------------------------------------------------------

\begin{align*}
  R_{0} &\leftarrow \text{initial orientation estimate
    (e.g., from accelerometer + magnetometer)}, \\
  \mathbf{b}_{0} &\leftarrow \mathbf{0}, \\
  P_{0} &\leftarrow \text{diag}\bigl(\sigma_{\phi,0}^{2}\,\eye_{3},\;
    \sigma_{b,0}^{2}\,\eye_{3}\bigr)
    \in \R^{6 \times 6}, \\
  \delta\mathbf{x}_{0} &\leftarrow \mathbf{0} \in \R^{6}.
\end{align*}

The error state is always reset to zero after each update (see
Step~3).

%----------------------------------------------------------------------
\subsection{Step 1: Prediction (Propagation)}
%----------------------------------------------------------------------

Given a gyroscope measurement $\bm{\omega}_m$ and time step
$\Delta t$:

\medskip
\noindent\textbf{1a.\ Propagate the state on the group:}
\begin{align}
  \bm{\omega} &= \bm{\omega}_m - \mathbf{b}, \label{eq:omega} \\
  R &\leftarrow R \;\expmap(\bm{\omega}\,\Delta t).
    \label{eq:Rprop}
\end{align}
(The bias estimate $\mathbf{b}$ is unchanged in the prediction
step.)

Equation~\eqref{eq:Rprop} is the \emph{group propagation}:
we integrate the rotation by multiplying on the right by a small
rotation $\expmap(\bm{\omega}\,\Delta t)$.  The result is
guaranteed to remain in $\SO$---no normalisation needed.

\medskip
\noindent\textbf{1b.\ Propagate the error-state covariance:}

The error-state dynamics (linearised) are:
\[
  \delta\dot{\mathbf{x}}
  = F\,\delta\mathbf{x} + G\,\mathbf{n},
\]
where $\mathbf{n} = (\mathbf{n}_{\omega}^{\mathsf{T}},\;
\mathbf{n}_{b}^{\mathsf{T}})^{\mathsf{T}}$ and:
\[
  F = \begin{pmatrix}
    -\bm{\omega}^{\wedge} & -\eye_{3} \\[3pt]
    \mathbf{0} & \mathbf{0}
  \end{pmatrix},
  \qquad
  G = \begin{pmatrix}
    -\eye_{3} & \mathbf{0} \\[3pt]
    \mathbf{0} & \eye_{3}
  \end{pmatrix}.
\]

\textbf{Where does $F$ come from?}  The upper-left block
$-\bm{\omega}^{\wedge}$ arises from linearising the group
propagation~\eqref{eq:Rprop} with respect to~$\delta\bm{\phi}$.
The upper-right block~$-\eye$ reflects the fact that a bias error
acts as a negative angular velocity error.  The lower row is zero
because the bias random walk has no state-dependent dynamics.

Discretise to get the state transition matrix and process noise:
\begin{align}
  F_{d} &= \exp(F\,\Delta t)
    \approx \eye_{6} + F\,\Delta t, \label{eq:Fd}\\
  Q_{d} &\approx G\,Q_{c}\,G^{\mathsf{T}}\,\Delta t,
    \label{eq:Qd}
\end{align}
where $Q_{c} = \text{diag}(\sigma_{\omega}^{2}\,\eye_{3},\;
\sigma_{b}^{2}\,\eye_{3})$.

The covariance prediction is then the standard EKF equation:
\[
  \boxed{P \leftarrow F_{d}\,P\,F_{d}^{\mathsf{T}} + Q_{d}.}
\]

\textbf{Note:}  If you want to use the exact right Jacobian
$\mathbf{J}_{r}(\bm{\omega}\,\Delta t)$ instead of the
first-order approximation in~\eqref{eq:Fd}, replace $F_d$ with:
\[
  F_{d} = \begin{pmatrix}
    \expmap(-\bm{\omega}\,\Delta t) &
    -\mathbf{J}_{r}(\bm{\omega}\,\Delta t)\,\Delta t \\
    \mathbf{0}_{3} & \eye_{3}
  \end{pmatrix}.
\]
For small $\Delta t$ (100\,Hz IMU), the first-order approximation
is usually sufficient.

%----------------------------------------------------------------------
\subsection{Step 2: Update (Correction)}
%----------------------------------------------------------------------

When an accelerometer measurement $\mathbf{a}_m$ arrives:

\medskip
\noindent\textbf{2a.\ Predicted measurement:}

The filter predicts what the accelerometer \emph{should} read given
the current orientation estimate:
\[
  \hat{\mathbf{a}} = R^{\mathsf{T}}\,\mathbf{g}.
\]

\noindent\textbf{2b.\ Innovation (residual):}
\[
  \mathbf{z} = \mathbf{a}_m - \hat{\mathbf{a}} \in \R^{3}.
\]
This is an ordinary vector subtraction in $\R^{3}$---no Lie theory
needed for the measurement residual because the accelerometer
output lives in $\R^{3}$.

\medskip
\noindent\textbf{2c.\ Measurement Jacobian:}

We need $H = \partial\mathbf{a}/\partial(\delta\mathbf{x})$, the
sensitivity of the predicted measurement to the error state.

A small orientation error $\delta\bm{\phi}$ perturbs the predicted
measurement as:
\begin{align*}
  \hat{\mathbf{a}}(\delta\bm{\phi})
  &= \bigl(\expmap(\delta\bm{\phi})\,R\bigr)^{\mathsf{T}}\,
    \mathbf{g}
  \\
  &= R^{\mathsf{T}}\,
    \expmap(\delta\bm{\phi})^{\mathsf{T}}\,\mathbf{g}
  \\
  &\approx R^{\mathsf{T}}\,
    \bigl(\eye - \delta\bm{\phi}^{\wedge}\bigr)\,\mathbf{g}
  \\
  &= R^{\mathsf{T}}\,\mathbf{g}
    + R^{\mathsf{T}}\,\mathbf{g}^{\wedge}\,\delta\bm{\phi}.
\end{align*}
(We used
$\expmap(\delta\bm{\phi})^{\mathsf{T}}
\approx \eye - \delta\bm{\phi}^{\wedge}$
and the cross-product identity
$-\mathbf{a}^{\wedge}\mathbf{b}
= \mathbf{b}^{\wedge}\mathbf{a}$.)

Therefore:
\[
  \boxed{
  H = \begin{pmatrix}
    R^{\mathsf{T}}\,\mathbf{g}^{\wedge}
    & \mathbf{0}_{3 \times 3}
  \end{pmatrix}
  \in \R^{3 \times 6}.
  }
\]
The bias columns are zero because the accelerometer measurement
does not depend on the gyro bias.

\medskip
\noindent\textbf{2d.\ Standard EKF update:}

With the measurement noise covariance
$V = \sigma_{a}^{2}\,\eye_{3}$:
\begin{align}
  S &= H\,P\,H^{\mathsf{T}} + V, \label{eq:S}\\
  K &= P\,H^{\mathsf{T}}\,S^{-1}, \label{eq:K}\\
  \delta\mathbf{x} &= K\,\mathbf{z}, \label{eq:dx}\\
  P &\leftarrow (\eye_{6} - K\,H)\,P. \label{eq:Pupd}
\end{align}

Equations~\eqref{eq:S}--\eqref{eq:Pupd} are \emph{exactly} the
standard EKF update equations.  Nothing has changed.  The Lie group
structure is entirely hidden in how we defined $H$ and what we do
with $\delta\mathbf{x}$ next.

%----------------------------------------------------------------------
\subsection{Step 3: Injection (Reset)}
%----------------------------------------------------------------------

This is the step that makes the filter a ``Lie group filter.''  The
Kalman update produced a correction
$\delta\mathbf{x} = (\delta\bm{\phi}^{\mathsf{T}},\;
\delta\mathbf{b}^{\mathsf{T}})^{\mathsf{T}}$.  We inject it into
the state:

\[
  \boxed{
  \begin{aligned}
    R &\leftarrow \expmap(\delta\bm{\phi})\; R, \\
    \mathbf{b} &\leftarrow \mathbf{b} + \delta\mathbf{b}, \\
    \delta\mathbf{x} &\leftarrow \mathbf{0}.
  \end{aligned}
  }
\]

\textbf{What is happening:}
\begin{itemize}[nosep]
  \item The orientation correction $\delta\bm{\phi}$ is mapped from
    $\R^{3}$ to $\SO$ via the exponential map, then composed with
    the current estimate by \emph{matrix multiplication} (not
    addition).
  \item The bias correction is ordinary vector addition.
  \item The error state is reset to zero, because the correction has
    been absorbed into the ``nominal'' state.
\end{itemize}

\textbf{This is the entire Lie group content of the filter.}  The
exponential map is used exactly once per update cycle, in a single
line of code.  Everything else is standard EKF.

\newpage
%======================================================================
\section{Complete Algorithm}
%======================================================================

\begin{algorithm}[H]
\caption{SO(3) Error-State EKF for Attitude Estimation}
\begin{algorithmic}[1]
\State \textbf{Initialise:} $R$, $\mathbf{b} \leftarrow \mathbf{0}$,
  $P \leftarrow P_0$
\For{each time step}
  \State \textbf{--- Prediction ---}
  \State $\bm{\omega} \leftarrow \bm{\omega}_m - \mathbf{b}$
    \Comment{Subtract bias}
  \State $R \leftarrow R \;\expmap(\bm{\omega}\,\Delta t)$
    \Comment{Propagate on group}
  \State $F_{d} \leftarrow \eye_{6}
    + \begin{psmallmatrix}
      -\bm{\omega}^{\wedge} & -\eye_{3} \\
      \mathbf{0} & \mathbf{0}
    \end{psmallmatrix} \Delta t$
    \Comment{Error-state Jacobian}
  \State $Q_{d} \leftarrow G\,Q_{c}\,G^{\mathsf{T}}\,\Delta t$
  \State $P \leftarrow F_{d}\,P\,F_{d}^{\mathsf{T}} + Q_{d}$
    \Comment{Covariance prediction}
  \Statex
  \If{accelerometer measurement available}
    \State \textbf{--- Update ---}
    \State $\hat{\mathbf{a}} \leftarrow R^{\mathsf{T}}\,\mathbf{g}$
      \Comment{Predicted measurement}
    \State $\mathbf{z} \leftarrow \mathbf{a}_m - \hat{\mathbf{a}}$
      \Comment{Innovation}
    \State $H \leftarrow
      \begin{psmallmatrix}
        R^{\mathsf{T}}\mathbf{g}^{\wedge} & \mathbf{0}_{3\times 3}
      \end{psmallmatrix}$
      \Comment{Measurement Jacobian}
    \State $S \leftarrow H\,P\,H^{\mathsf{T}} + V$
    \State $K \leftarrow P\,H^{\mathsf{T}}\,S^{-1}$
      \Comment{Kalman gain}
    \State $\delta\mathbf{x} \leftarrow K\,\mathbf{z}$
      \Comment{Correction}
    \State $P \leftarrow (\eye_{6} - K\,H)\,P$
      \Comment{Covariance update}
    \Statex
    \State \textbf{--- Injection ---}
    \State $R \leftarrow
      \expmap(\delta\mathbf{x}_{1:3})\; R$
      \Comment{Apply rotation correction}
    \State $\mathbf{b} \leftarrow
      \mathbf{b} + \delta\mathbf{x}_{4:6}$
      \Comment{Apply bias correction}
  \EndIf
\EndFor
\end{algorithmic}
\end{algorithm}

\textbf{That's it.}  The entire filter is 15 lines of mathematics.
The only function you need beyond standard linear algebra is
$\expmap(\cdot)$, which is the three-line Rodrigues formula.

\newpage
%======================================================================
\section{Implementation Notes}
%======================================================================

\subsection{The Rodrigues Formula in Code}

The exponential map is the only Lie-specific function you need to
implement.  In pseudocode:

\begin{verbatim}
function R = Exp_SO3(phi)
    theta = norm(phi);
    if theta < 1e-10
        R = eye(3) + skew(phi);   % first-order approx
    else
        a = phi / theta;          % unit axis
        K = skew(a);              % skew-symmetric matrix
        R = eye(3) + sin(theta)*K + (1-cos(theta))*K^2;
    end
end
\end{verbatim}

And the skew-symmetric matrix:

\begin{verbatim}
function S = skew(v)
    S = [  0   -v(3)  v(2);
          v(3)   0   -v(1);
         -v(2)  v(1)   0  ];
end
\end{verbatim}

\subsection{Numerical Considerations}

\begin{itemize}
  \item \textbf{Orthogonalisation:}  Due to floating-point
    arithmetic, $R$ will slowly drift away from $\SO$ over many
    propagation steps.  Periodically re-orthogonalise using an SVD:
    $R \leftarrow U\,V^{\mathsf{T}}$ where $U\Sigma V^{\mathsf{T}}
    = \text{SVD}(R)$.  This is needed perhaps every 1000 steps.
  \item \textbf{Small-angle guard:}  In $\expmap$ and $\logmap$,
    use Taylor expansions when $\theta < 10^{-10}$ to avoid
    division by zero.
  \item \textbf{Symmetry of $P$:}  After the covariance update,
    enforce symmetry: $P \leftarrow (P + P^{\mathsf{T}})/2$.
  \item \textbf{IMU rate vs.\ accelerometer rate:}  Gyroscopes
    often run at 100--1000\,Hz while accelerometers may be slower.
    Run the prediction step at the gyro rate; run the update step
    only when accelerometer data arrives.
\end{itemize}

\subsection{Tuning Parameters}

\begin{center}
\begin{tabular}{lll}
\toprule
\textbf{Parameter} & \textbf{Typical value} & \textbf{Meaning} \\
\midrule
$\sigma_{\omega}$ & $0.01$--$0.1$ rad/s & Gyro noise density \\
$\sigma_{b}$ & $10^{-4}$--$10^{-3}$ rad/s$^{2}$ & Bias random walk \\
$\sigma_{a}$ & $0.5$--$2.0$ m/s$^{2}$ & Accelerometer noise \\
$\sigma_{\phi,0}$ & $0.1$--$1.0$ rad & Initial orientation uncertainty \\
$\sigma_{b,0}$ & $0.01$--$0.1$ rad/s & Initial bias uncertainty \\
\bottomrule
\end{tabular}
\end{center}

These values come from the IMU datasheet (noise spectral density
and bias instability) or from empirical tuning.

\newpage
%======================================================================
\section{What Makes This a ``Lie Group Filter''?}
%======================================================================

Let us step back and see what we have done.

\subsection{Comparison with the Standard Vector-State EKF}

\begin{center}
\begin{tabular}{lll}
\toprule
& \textbf{Euler-angle EKF} & \textbf{SO(3) error-state EKF} \\
\midrule
State & $(\phi, \theta, \psi) \in \R^{3}$
  & $R \in \SO$ \\
Error state & $\delta(\phi, \theta, \psi) \in \R^{3}$
  & $\delta\bm{\phi} \in \R^{3} \cong \so$ \\
Propagation & Euler angle kinematics & $R \leftarrow R\,\expmap(\bm{\omega}\,\Delta t)$ \\
Correction & $\bm{\alpha} \leftarrow \bm{\alpha} + \delta\bm{\alpha}$
  & $R \leftarrow \expmap(\delta\bm{\phi})\,R$ \\
Singularity? & Yes (gimbal lock) & No \\
Covariance dim. & $3 \times 3$ (or $6 \times 6$)
  & $6 \times 6$ (same) \\
\bottomrule
\end{tabular}
\end{center}

The filter structure is identical.  The only differences are:
\begin{enumerate}[nosep]
  \item The state $R$ is propagated using matrix multiplication
    instead of integrating Euler angle rates.
  \item The correction is applied via the exponential map (one
    matrix multiplication) instead of vector addition.
  \item There is no gimbal lock, no normalisation constraint, and
    no singularity---ever.
\end{enumerate}

\subsection{The Pattern Generalises}

The structure you just learned---\emph{propagate on the group,
correct in the algebra}---is the same for every Lie group filter:

\begin{center}
\begin{tabular}{ll}
\toprule
\textbf{Application} & \textbf{Lie group} \\
\midrule
Attitude estimation & $\SO$ \\
Pose estimation (position + orientation) & $\mathrm{SE}(3)$ \\
Inertial navigation (pose + velocity) & $\mathrm{SE}_{2}(3)$ \\
2D robot localisation & $\mathrm{SE}(2)$ \\
\bottomrule
\end{tabular}
\end{center}

In each case, you need only the same five tools (group, algebra,
$\expmap$, $\logmap$, Jacobians) specialised to the group in
question.  The Sol\`{a} ``Micro Lie Theory'' paper provides the
formulas for all of these groups in a single reference.

\newpage
%======================================================================
\section{Where to Go Next}
%======================================================================

\begin{enumerate}
  \item \textbf{Implement it.}  Build the SO(3) filter in
    MATLAB or Python.  Generate synthetic gyro and accelerometer
    data from a known rotation trajectory, run the filter, and
    compare to ground truth.  The most clarifying exercise.

  \item \textbf{Read Sol\`{a}'s ``Micro Lie Theory'' paper}
    (\href{https://arxiv.org/abs/1812.01537}{arXiv:1812.01537}).
    You now have the vocabulary to read it fluently.  It generalises
    everything in this document to arbitrary Lie groups and provides
    the formulas for SE(3) and SE$_{2}$(3).

  \item \textbf{Read Sol\`{a}'s ``Quaternion Kinematics for the
    ESKF''}
    (\href{https://arxiv.org/abs/1711.02508}{arXiv:1711.02508}).
    This gives a detailed treatment of the same filter using
    quaternions internally (but with the same Lie algebra error
    state).

  \item \textbf{Study Barfoot's \emph{State Estimation for
    Robotics}} (free draft at
    \href{http://asrl.utias.utoronto.ca/~tdb/bib/barfoot_ser24.pdf}%
    {asrl.utias.utoronto.ca}).  Chapters on rotations and poses,
    then the estimation chapters.

  \item \textbf{Learn the Invariant EKF.}  Read Barrau and
    Bonnabel's Annual Review paper
    (\href{https://hal.science/hal-05088150v1/file/AR_barrau_bonnabel.pdf}%
    {HAL preprint}).  The IEKF uses a different error
    convention (left- vs.\ right-invariant) that yields
    trajectory-independent error dynamics and provable convergence
    guarantees.

  \item \textbf{Explore the Equivariant Filter.}  van Goor, Hamel,
    and Mahony's work
    (\href{https://arxiv.org/abs/2010.14666}{arXiv:2010.14666})
    represents the current frontier of geometric filtering.
\end{enumerate}

\subsection*{Key References}

\begin{itemize}[nosep,leftmargin=*]
  \item Sol\`{a} et al.\ ``A Micro Lie Theory for State
    Estimation in Robotics.'' arXiv:1812.01537.
  \item Sol\`{a}, ``Quaternion Kinematics for the Error-State
    Kalman Filter,'' arXiv:1711.02508, 2017.
  \item Barfoot, \emph{State Estimation for Robotics,} 2nd~ed.,
    Cambridge, 2024.
  \item Barrau and Bonnabel, ``Invariant Kalman Filtering,''
    \emph{Annual Review of Control, Robotics, and Autonomous
    Systems,} 2018.
  \item Mahony, Hamel, Pflimlin, ``Nonlinear Complementary Filters
    on the Special Orthogonal Group,'' \emph{IEEE Trans.\ Automatic
    Control,} 2008.
\end{itemize}

\end{document}
